\documentclass[12pt,a4paper]{article}

\usepackage[utf8]{inputenc}
\usepackage[spanish]{babel}
\usepackage{amsmath, amssymb}
\usepackage{graphicx}
\usepackage{float}
\usepackage{tabularx}
\usepackage{geometry}
\usepackage{enumitem}
\usepackage{titlesec}
\usepackage{booktabs}
\usepackage{array}

\geometry{margin=2.5cm}
\titleformat{\section}{\large\bfseries}{\thesection.}{1em}{}

\begin{document}

\begin{center}
    \Large \textbf{Taller 3 - Diagnóstico de Viabilidad Empírica del Proyecto} \\
    \vspace{0.3cm}
    \large \textbf{Proyecto:} \textit{Estimación de Probabilidad de Victoria en CS:GO Pre-match} \\
    \vspace{0.5cm}
    \normalsize \textbf{Integrantes:} Julian Jimenez, Julian Duarte, Tomas Rincon \\
    \textbf{Fecha:} \today
\end{center}

\vspace{0.8cm}

\section{1. Matriz de Viabilidad Completa}

A continuación se presenta la matriz detallada con las variables fundamentales para la estimación bayesiana de probabilidad de victoria en encuentros de élite.

\begin{table}[H]
\centering
\resizebox{\textwidth}{!}{%
\begin{tabular}{|l|p{3cm}|p{3cm}|p{3cm}|p{3cm}|p{3cm}|}
\hline
\textbf{Elemento} & \textbf{Definición actual} & \textbf{Variable propuesta} & \textbf{Fuente de datos} & \textbf{Nivel (Unidad)} & \textbf{Problema identificado} \\ \hline
\textbf{Var. Dependiente} & Resultado binario de la partida. & \texttt{match\_winner\_1} (Boo: 1 si Team 1 gana, 0 si pierde). & HLTV Dataset (Kaggle: "CS:GO Professional Matches"). & Partida (Match). & Ninguno; el dato es objetivo y etiquetado nativamente. \\ \hline
\textbf{Var. Indep. 1} & Ranking relativo de los equipos. & \texttt{ranking\_delta} (Dif. logarítmica de puestos en el Top 30). & HLTV Rankings (Scraper/Snapshot histórico). & Equipo (Team ID). & El ranking es un indicador retardado (no captura cambios de roster inmediatos). \\ \hline
\textbf{Var. Indep. 2} & Desempeño individual agregado. & \texttt{avg\_player\_rating\_3m} (Promedio de Rating 2.0 de los 5 jugadores). & Player Stats (HLTV - 3 months leading to match). & Jugador (agregado por quinteto). & Difícil de rastrear cuando hay sustitutos temporales (\textit{stand-ins}). \\ \hline
\textbf{Var. Indep. 3} & Historial en el mapa específico. & \texttt{map\_winrate\_6m} (Diferencia de \% de victorias en el mapa a jugar). & Match history (HLTV filtered by Map). & Mapa por Equipo. & Colinealidad con el nivel general del equipo; escasez de datos en mapas nuevos (ej. Anubis). \\ \hline
\end{tabular}%
}
\caption{Matriz de viabilidad empírica para el modelo predictivo de CS:GO.}
\end{table}

\section{2. Diagnóstico de Diseño (Parte Proceso)}

Como parte del ejercicio de diagnóstico, se responden las siguientes interrogantes clave sobre el diseño actual:

\begin{enumerate}
    \item \textbf{¿Cuál es el mayor problema de su diseño actual?}
    El mayor desafío radica en la \textbf{estacionariedad de los datos}. En CS:GO, el rendimiento de los equipos cambia drásticamente con cada actualización del parche; predecir el futuro basándose en un historial estático puede estar sesgado por una realidad competitiva ya obsoleta.
    
    \item \textbf{¿Qué variable es más débil?}
    La variable de **desempeño en el mapa**. Aunque es crucial tácticamente, es estadísticamente ruidosa debido a que el historial puede ser limitado para ciertos mapas poco jugados, distorsionando el poder real del equipo.
    
    \item \textbf{¿Qué dato aún no tienen?}
    Aún no se ha consolidado un \textbf{Pipeline de unión (Join)} automatizado que vincule por fecha cada partida con el ranking histórico exacto de ese día.
    
    \item \textbf{¿Qué parte del proyecto probablemente deben cambiar?}
    Es necesario restringir el alcance del modelo únicamente al \textbf{Tier 1 Profesional}, ya que la inconsistencia de datos en equipos de ligas menores impide la convergencia robusta del modelo Bayesiano.
\end{enumerate}

\section{3. Diagnóstico Crítico (Entregable Final)}

\subsection{Principales limitaciones de datos}
La limitación terminal es la volatilidad de los \textit{rosters}. Además, datos cualitativos vitales como el estado psicológico de los jugadores o la comunicación interna no están disponibles, limitando el modelo a factores puramente mecánicos.

\subsection{Decisiones tomadas}
Se ha decidido emplear un modelo de \textbf{Regresión Logística Bayesiana} para obtener una \textbf{cuantificación de la incertidumbre} (HDI), permitiendo que el modelo admita su propio desconocimiento frente a datos contradictorios.

\subsection{Cambios en el proyecto}
Se opta por \textbf{acotar el alcance} a equipos Top 30 mundial. Esto garantiza que la fuente de datos sea de alta fidelidad y que los jugadores operen en su máximo potencial, minimizando los factores aleatorios.

\section{4. Versión Ajustada del Proyecto}

\subsection{Pregunta de Investigación Refinada}
\textit{¿En qué medida el desempeño individual acumulado (Rating 2.0) y la jerarquía histórica (Ranking HLTV) permiten cuantificar la incertidumbre de victoria en encuentros de élite de CS:GO antes de su inicio?}

\subsection{Breve justificación del cambio}
El ajuste centra el valor del proyecto en la \textbf{calibración probabilística}. No buscamos simplemente una predicción binaria, sino entender qué tan seguros podemos estar de un resultado dado el historial de élite.

\end{document}
